% English translation of prelude.tex, lines 127--153.
% Source: Methods of Algebra, Volume 2, commit
% 9a5803ff77dd3257484cb177f851a73770a59dd3 (CC BY 4.0).
% stable-unit-id: o014.aljabr2.prelude.reader-guide

% segment-id: o014.li2.prelude.reader-guide.h001
\section*{Guide to Using the Book}
\addcontentsline{toc}{section}{Guide to Using the Book}
\hypertarget{o014.aljabr2.prelude.reader-guide}{ }

% segment-id: o014.li2.prelude.reader-guide.h002
\subsection*{How the Book May Be Used}

% segment-id: o014.li2.prelude.reader-guide.p001
To serve as a reference work, the book needs a reasonably systematic organization. Its chapters therefore follow a primarily logical order. Dependencies between earlier and later material are especially pronounced in the Inner Part, while the few necessary topics that diverge from the main line or are more technical have been assigned to appendices.

% segment-id: o014.li2.prelude.reader-guide.p002
Independent readers should decide what to read according to their own interests and needs. Most of these choices arise in the Outer Part, whereas the Inner Part supplies the common foundational framework, although even there some sections are optional. More abstract or technical details may be omitted on an initial reading, while material involving the appendices is best skipped at first or merely skimmed. Detailed advice on these choices appears at the beginning of each chapter and appendix, especially in the reading guides.

% segment-id: o014.li2.prelude.reader-guide.p003
As noted above in the discussion of the book's aims, we have unavoidably wrapped mathematical ideas in layer upon layer of prose. Independent readers must make a determined effort to discern the mathematics itself within those layers, like catching lightning amid roiling clouds.

% segment-id: o014.li2.prelude.reader-guide.p004
For readers using the book as a reference, the author has tried to organize the exposition modularly so that its parts can be read separately. This is reflected not only in the division into chapters, but also in the following two features.
\begin{itemize}
	% segment-id: o014.li2.prelude.reader-guide.i001
	\item As far as possible, the book uses the standard notation of contemporary literature, supplemented by reviews and explanations. The Conventions at the end of the introduction and the index at the end of the book provide two further aids to understanding the notation.
	% segment-id: o014.li2.prelude.reader-guide.i002
	\item Statements and proofs contain copious cross-references, allowing readers who already know something of the subject to understand the discussion quickly and build an efficient conceptual map.
\end{itemize}

% segment-id: o014.li2.prelude.reader-guide.p005
In particular, the author hopes that the book will help readers preparing to teach or participating in seminars to assemble the knowledge they need quickly and then recast it in their own language.

% segment-id: o014.li2.prelude.reader-guide.p006
As for classroom use, one must frankly admit that teaching schedules are not governed by individual wishes. While writing the book, the author had almost no opportunity to try out this material in the classroom, so its effectiveness there remains unknown. Judging from past experience, even with details omitted in lectures, the book clearly contains enough material for more than one academic year; in a streamlined presentation, the Inner Part might perhaps be compressed into a single semester. An instructor using the book as the primary textbook will first have to decide how to divide the material. The author hopes that the modular design just mentioned will make this task easier.

% segment-id: o014.li2.prelude.reader-guide.p007
The lack of classroom testing has two further consequences. First, the book must still contain many undiscovered errors. Second, it lacks the finish of a classic textbook. This is due both to practical circumstances and to the author's own limitations; the author can only ask the reader's indulgence.

% segment-id: o014.li2.prelude.reader-guide.h003
\subsection*{On the Exercises}

% segment-id: o014.li2.prelude.reader-guide.p008
Every chapter and appendix ends with exercises arranged in the order of the corresponding material in the main text and sometimes accompanied by hints. Their chief purposes are to supplement the text with further properties, provide examples, and broaden the reader's perspective. Some exercises help readers become familiar with techniques and generally make few technical demands. A few ask readers to fill in details omitted from the text; these details are either trivial, easy, or, more often, both. In short, the author hopes that beginners will attempt as many exercises as possible, but they need not force themselves or feel obliged to complete them all.

% segment-id: o014.li2.prelude.reader-guide.h004
\subsection*{Categories and Their Size}

% segment-id: o014.li2.prelude.reader-guide.p009
Most of this book is expressed in categorical language because practice has shown category theory to be an efficient means of understanding and handling algebraic problems. Nevertheless, linear algebra---or what is called homological algebra---is no more a branch of category theory than probability theory is a branch of measure theory.

% segment-id: o014.li2.prelude.reader-guide.p010
Unlike some textbooks, both this book and Volume~I require that, for a category, the collection of all objects and the collection of all morphisms each be a set; it is not enough for the objects merely to form a ``class.'' The price of excluding proper classes from the definitions is that, strictly speaking, one must introduce Grothendieck universes to measure the size of sets and assume that all sets belong to some universe $\mathcal{U}$; see \cite[Assumption~1.5.2, Remark~1.5.6]{Li1}. When concrete problems are considered case by case, the axiom that there are sufficiently many universes---equivalently, sufficiently many strongly inaccessible cardinals---can often be weakened. Since this issue is not central to the book, we shall not pursue it further.
